% Pursuing Stacks (cotes 134-2 to 134-8): archive cote 134-4, site page 34 (PDF page 35), typescript page 122.
% typescript. Draft compared: ps3.tex:1677-1752.
% AI-transcribed by Claude Opus 5.5 (claude-opus-5-5, adaptive thinking, effort high), Message Batches API, batch ps2 (msgbatch_01NVwr1Yuo8qJf5rykeGBXPK), 2026-11-12; not yet checked by eye.
\scanpage{134-4/34}
\mspage{122}
Moreover, \xout{if} this condition \xout{\unsure{is satisfied, the}}\inserted[typed]{implies the} following two:

c)\quad For any object Y in M, if \underline{Hom}(Y,X) exists, it is contractible.

d)\quad For any Y in M, if Hom(X,Y) exists, the natural map
\[ Y \hand{\to} \underline{Hom}(X,Y) \]
(identifying Y to the ``subobject of constant maps from X to Y'') is a
homotopism.

The equivalence of a) and b) is just a tautological translation
of \xout{\unsure{contractibility}}\inserted[typed]{contractibility} \inserted{and homotopy relation} in terms of \deleted{\unsure{elementary}} weak homotopy intervals $\underline{J}$\inserted{(cf. cor. 2 below)}.
T\hand{h}at c\hand{)} and d\hand{)} follow comes from the fact that the monoïd\xout{\unsure{e}} object
$I = \underline{End}(X,X) = \underline{Hom}(X,X)$ operates on the left on \underline{Hom}(Y,X), on the right
on \underline{Hom}(X,Y), and the following :

\uline{Corollary}\hand{1}. Let $\underline{I}$ be a \xout{\unsure{contractible contractor}}\inserted[typed]{weak homotopy interval}, assume the underlying I
``operates'' on an object H , namely we are given a map
\[ h: I\mathrm{x}H \hand{\longrightarrow} H \qquad (\text{``operation'' of I on H}) \]
satisfying the relations (where h(u,f) is written simply u.f ):
\[ \xout{xx}\ e_o.f = f \quad , \quad e_1\hand{(}e_1\hand{.}f) = e_1.\ f \quad , \]
namely $e_o$ acts as the identity, and $\hand{e_1}$ acts as an idempotent \inserted[typed]{p on H} (a very
weak associativity assumption indeed\inserted[typed]{if I is a contractor}, as $e_1\hand{\unsure{\circ|\cdot}}e_1 = e_1$). Assume the image
of p i.e. Ker($id_H$,p) exists, let $H_o$ be the corresponding subobject of
H, and \deleted{p}
\[ p_o : H \hand{\longrightarrow} H_o \]
the map induced by p. Then $p_o$ is a homotopism (and hence the inclusion
\inserted[typed]{i:}$H_o \hand{\longrightarrow} H$, which is a section of $p_o$ , is a homotopism too).

Because of the saturation \xout{condition}\inserted[typed]{property c')} of homotopies, it is enough to
check that $p = ip_o$ is a homotopism (as $pi_o = id_{H_o}$ already is one), and
for this it is enough to see it is homotopic to the identity map of H.
But a homotopy between the two is realized by h, \deleted{\unsure{hence}} qed.

The argument for equivalence of a) and b) above can be generalized
as follows:

\uline{Corollary} 2 . Let M be as before, and X and Y objects such that
$H = \underline{Hom}(X,Y)$ exists in M. Let $f,g: X \hand{\rightrightarrows} Y$ be two maps, which we'll
identify to the corresponding sections of H. Then f and g are homotopic
maps iff they gi\hand{v}e rise to homotopic sections of H.

F)\quad \uline{The canonical homotopy structure} : \hand{\uline{preliminaries on} $\pi_o$ .}

In order to simplify life, I will in this section make the follo-%
wing assumptions on M (which presumably, except for the first, could be

\margin{34}
